\chapter{Cơ sở Lý thuyết}

\section{Mạng Nơ-ron Tích chập Đồ thị (GCN)}

\subsection{Lý thuyết Đồ thị và Biểu diễn Mạng lưới Giao thông}

Mạng lưới đường phố được mô hình hóa tự nhiên bằng \textbf{đồ thị có hướng} $\mathcal{G} = (\mathcal{V}, \mathcal{E}, \mathbf{W})$, trong đó:

\begin{itemize}[itemsep=3pt]
  \item $\mathcal{V} = \{v_1, v_2, \ldots, v_N\}$: Tập hợp $N$ \textbf{đỉnh} (nodes) --- mỗi đỉnh đại diện cho một nút giao hoặc đoạn đường có gắn camera.
  \item $\mathcal{E} \subseteq \mathcal{V} \times \mathcal{V}$: Tập hợp các \textbf{cạnh} (edges) --- mỗi cạnh $(v_i, v_j)$ biểu diễn đoạn đường nối trực tiếp từ nút $v_i$ đến nút $v_j$.
  \item $\mathbf{W} \in \mathbb{R}^{N \times N}$: \textbf{Ma trận trọng số} --- $W_{ij}$ phản ánh mức độ ảnh hưởng giao thông từ nút $i$ đến nút $j$, thường được tính bằng hàm giảm theo khoảng cách hoặc thời gian di chuyển.
\end{itemize}

Trọng số thường được thiết kế theo công thức Gaussian:
\begin{equation}
  W_{ij} = \exp\!\left(-\frac{d_{ij}^2}{\sigma^2}\right)
  \label{eq:edge_weight}
\end{equation}
với $d_{ij}$ là khoảng cách (km) hoặc thời gian di chuyển (phút) giữa hai nút, $\sigma$ là tham số điều chỉnh độ rộng.

\subsection{Phép Tích chập Trên Đồ thị}

Khác với ảnh 2D có lưới đều đặn, dữ liệu trên đồ thị không có cấu trúc lưới. Phép tích chập đồ thị (Graph Convolution) được định nghĩa trong miền phổ (spectral domain) thông qua phân tích riêng (eigen-decomposition) của ma trận Laplacian chuẩn hóa:

\begin{equation}
  \mathbf{L} = \mathbf{I}_N - \mathbf{D}^{-1/2}\mathbf{A}\mathbf{D}^{-1/2}
  \label{eq:laplacian}
\end{equation}

Kipf \& Welling \cite{kipf2017gcn} đề xuất xấp xỉ bậc nhất (first-order approximation) rất hiệu quả:

\begin{equation}
  \mathbf{H}^{(l+1)} = \sigma\!\left(\tilde{\mathbf{D}}^{-1/2}\tilde{\mathbf{A}}\tilde{\mathbf{D}}^{-1/2}\mathbf{H}^{(l)}\mathbf{W}^{(l)}\right)
  \label{eq:gcn_update}
\end{equation}

trong đó:
\begin{itemize}[itemsep=3pt]
  \item $\tilde{\mathbf{A}} = \mathbf{A} + \mathbf{I}_N$: Ma trận kề có thêm self-loop (để mỗi nút nhận thông tin từ chính nó).
  \item $\tilde{\mathbf{D}}_{ii} = \sum_j \tilde{A}_{ij}$: Ma trận bậc (degree matrix) của $\tilde{\mathbf{A}}$.
  \item $\mathbf{H}^{(l)} \in \mathbb{R}^{N \times F_l}$: Ma trận đặc trưng ẩn tại lớp $l$, với $F_l$ chiều đặc trưng.
  \item $\mathbf{W}^{(l)} \in \mathbb{R}^{F_l \times F_{l+1}}$: Ma trận trọng số học được của lớp $l$.
  \item $\sigma(\cdot)$: Hàm kích hoạt phi tuyến (ReLU, LeakyReLU).
\end{itemize}

\begin{tcolorbox}[tipbox]
\textbf{Ý nghĩa vật lý:} Phép tích chập đồ thị cho phép mô hình học cách \textit{ùn tắc lan truyền}: nếu nút giao A bắt đầu tắc, sau một lớp GCN, thông tin này được truyền sang các nút lân cận B, C; sau hai lớp, thông tin đạt đến các nút bậc 2 (lân cận của lân cận).
\end{tcolorbox}

\section{Mạng LSTM và Biến thể Stacked}

\subsection{Cơ chế Gate của LSTM}

Mạng LSTM \cite{hochreiter1997lstm} giải quyết vấn đề Vanishing Gradient của RNN thông thường thông qua cơ chế \textbf{cổng (gate)} điều tiết luồng thông tin. Tại mỗi bước thời gian $t$:

\begin{align}
  \mathbf{f}_t &= \sigma(\mathbf{W}_f[\mathbf{h}_{t-1}, \mathbf{x}_t] + \mathbf{b}_f) \quad \text{(Forget Gate)} \label{eq:forget} \\
  \mathbf{i}_t &= \sigma(\mathbf{W}_i[\mathbf{h}_{t-1}, \mathbf{x}_t] + \mathbf{b}_i) \quad \text{(Input Gate)} \label{eq:input} \\
  \tilde{\mathbf{c}}_t &= \tanh(\mathbf{W}_c[\mathbf{h}_{t-1}, \mathbf{x}_t] + \mathbf{b}_c) \quad \text{(Candidate Cell)} \label{eq:candidate} \\
  \mathbf{c}_t &= \mathbf{f}_t \odot \mathbf{c}_{t-1} + \mathbf{i}_t \odot \tilde{\mathbf{c}}_t \quad \text{(Cell State Update)} \label{eq:cell} \\
  \mathbf{o}_t &= \sigma(\mathbf{W}_o[\mathbf{h}_{t-1}, \mathbf{x}_t] + \mathbf{b}_o) \quad \text{(Output Gate)} \label{eq:output_gate} \\
  \mathbf{h}_t &= \mathbf{o}_t \odot \tanh(\mathbf{c}_t) \quad \text{(Hidden State)} \label{eq:hidden}
\end{align}

trong đó $\odot$ là tích Hadamard (element-wise multiplication), $\sigma$ là hàm sigmoid, và $[\cdot, \cdot]$ là phép ghép nối (concatenation).

\subsection{Stacked LSTM và Khả năng Biểu diễn Bậc cao}

Khi xếp chồng nhiều lớp LSTM (Stacked LSTM), mỗi lớp học một cấp độ trừu tượng khác nhau của chuỗi thời gian:

\begin{table}[H]
\centering
\caption{Ý nghĩa phân cấp trong Stacked LSTM áp dụng cho giao thông}
\label{tab:stacked_lstm}
\renewcommand{\arraystretch}{1.4}
\begin{tabular}{>{\bfseries\centering}p{1.8cm} p{4cm} p{6.5cm}}
\toprule
\textbf{Lớp} & \textbf{Đặc trưng học được} & \textbf{Ví dụ cụ thể} \\
\midrule
LSTM-1 & Biến động ngắn hạn (5--15 phút) & Đột biến lưu lượng do đèn đỏ \\
\rowcolor{rowgray}
LSTM-2 & Xu hướng trung hạn (15--45 phút) & Sóng ùn tắc từ một sự cố lan ra \\
LSTM-3 & Chu kỳ giờ cao điểm (1--2 giờ) & Sáng đông xe -- Trưa thưa -- Chiều kẹt \\
\bottomrule
\end{tabular}
\end{table}

\section{Mô hình Thay thế Nơ-ron (Neural Surrogate Model)}

\subsection{Hạn chế của Mô phỏng Vi mô}

Các công cụ mô phỏng vi mô (Microscopic Simulation) như SUMO (Simulation of Urban MObility) hay Vissim cho phép mô phỏng từng phương tiện riêng lẻ với độ chính xác vật lý cao. Tuy nhiên, độ phức tạp tính toán là $\mathcal{O}(N^2)$ với $N$ là số phương tiện:

\begin{table}[H]
\centering
\caption{So sánh Mô phỏng Vi mô và Surrogate Model}
\label{tab:sim_comparison}
\renewcommand{\arraystretch}{1.4}
\begin{tabular}{l p{4cm} p{4cm}}
\toprule
\textbf{Tiêu chí} & \textbf{Mô phỏng Vi mô (SUMO)} & \textbf{Neural Surrogate} \\
\midrule
Độ phức tạp & Phụ thuộc số xe, bước thời gian và mạng & Phụ thuộc số node, batch và kích thước model \\
\rowcolor{rowgray}
Thời gian chạy & 10 phút -- vài giờ & 100--300 ms \\
Độ chính xác & Rất cao (mô phỏng vật lý) & Xấp xỉ (học từ dữ liệu lịch sử) \\
\rowcolor{rowgray}
Xử lý Corner Cases & Phụ thuộc calibration và scenario coverage & Yêu cầu OOD gate và uncertainty calibration \\
Khả năng Real-time & Không & Có \\
\bottomrule
\end{tabular}
\end{table}

\subsection{Lý thuyết Deep Ensemble và ADE (Related Work)}

Deep Ensemble \cite{lakshminarayanan2017ade} là kỹ thuật đơn giản nhưng hiệu quả để ước lượng độ bất định (uncertainty):

\begin{equation}
  \hat{y}_{\text{ensemble}} = \frac{1}{M}\sum_{m=1}^{M} f_{\theta_m}(\mathbf{x}), \qquad
  \text{Var}[\hat{y}] = \frac{1}{M}\sum_{m=1}^{M}\left(f_{\theta_m}(\mathbf{x}) - \hat{y}_{\text{ensemble}}\right)^2
  \label{eq:ensemble}
\end{equation}

ADE (Adversarial Diverse Deep Ensemble) là hướng nghiên cứu liên quan. MVP STWI không triển khai ADE; hệ thống dùng heterogeneous ensemble và calibration trên held-out data.
\begin{enumerate}[itemsep=3pt]
  \item \textbf{Diversity:} Các sub-model $f_{\theta_1}, \ldots, f_{\theta_M}$ có kiến trúc khác nhau (MLP, CNN-1D, Transformer nhẹ) để tăng đa dạng dự báo.
  \item \textbf{Adversarial Training:} Huấn luyện với các mẫu đối kháng $\mathbf{x}' = \mathbf{x} + \epsilon \cdot \nabla_\mathbf{x} \mathcal{L}$ (theo FGSM) để mô hình bền vững với các kịch bản cực trị (Corner Cases).
\end{enumerate}

\section{Truy xuất Tăng cường Phát sinh (RAG)}

\subsection{Kiến trúc RAG Cổ điển}

RAG \cite{lewis2020rag} là kỹ thuật kết hợp LLM với nguồn tri thức ngoại vi để giảm thiểu ảo giác (hallucination). Quy trình gồm hai giai đoạn:

\begin{enumerate}[itemsep=4pt]
  \item \textbf{Retrieval (Truy xuất):} Cho câu truy vấn $q$, tính vector nhúng $\mathbf{e}_q = \text{Enc}(q)$, sau đó tìm $k$ đoạn văn bản $\{d_1, \ldots, d_k\}$ có độ tương đồng cosine cao nhất trong Vector Database:
    \begin{equation}
      \text{sim}(\mathbf{e}_q, \mathbf{e}_d) = \frac{\mathbf{e}_q \cdot \mathbf{e}_d}{\|\mathbf{e}_q\|\|\mathbf{e}_d\|}
      \label{eq:cosine}
    \end{equation}
  \item \textbf{Generation (Sinh văn bản):} LLM sinh câu trả lời dựa trên ngữ cảnh được bổ sung: $p(y|q, d_1, \ldots, d_k)$.
\end{enumerate}

\subsection{Schema-Level RAG và Text-to-SQL}

RAG truyền thống xử lý văn bản phi cấu trúc, còn số liệu mô phỏng cần truy vấn có kiểu. XiYanSQL \cite{guo2023xiyansql} là related work; MVP dùng LLM tạo \textit{SimulationQuery} có kiểu và server dựng SQL tham số hóa.

\section{Hệ thống Đa Tác tử (Multi-Agent Systems)}

\subsection{Định nghĩa và Phân loại Tác tử}

Một \textbf{tác tử} (agent) là thực thể tính toán có khả năng: \textit{nhận thức môi trường} (perception), \textit{suy luận} (reasoning) và \textit{hành động} (action) \cite{wooldridge2009multiagent}. Trong STWI, mỗi tác tử AI là một LLM được trang bị bộ công cụ (tools):

\begin{table}[H]
\centering
\caption{Phân loại và đặc điểm các Tác tử trong STWI}
\label{tab:agents}
\renewcommand{\arraystretch}{1.4}
\begin{tabular}{p{3.5cm} p{3.5cm} p{5.5cm}}
\toprule
\textbf{Tác tử} & \textbf{Công cụ (Tools)} & \textbf{Nhiệm vụ} \\
\midrule
Simulation Agent & API Tầng 2 (surrogate ensemble) & Chạy mô phỏng What-if \\
\rowcolor{rowgray}
Legal \& SOP Agent & Qdrant, SimulationQuery & Tra cứu pháp lý, truy vấn số liệu \\
Evaluation Agent & Tổng hợp kết quả & Đánh giá và đề xuất phương án \\
\rowcolor{rowgray}
Orchestrator & Điều phối 3 agent trên & Phân rã nhiệm vụ, hợp nhất kết quả \\
\bottomrule
\end{tabular}
\end{table}

\subsection{Kiến trúc ReAct và Chain-of-Thought}

Các tác tử trong STWI hoạt động theo kiến trúc \textbf{ReAct} (Reasoning + Acting): tại mỗi bước, tác tử luân phiên thực hiện \textit{Thought} (suy nghĩ), \textit{Action} (gọi tool), và \textit{Observation} (quan sát kết quả) cho đến khi đạt được mục tiêu.

\section{Suy luận Phản thực tế (Counterfactual Reasoning)}

\subsection{Lý thuyết Phản thực tế của Pearl}

Suy luận nhân quả (causal reasoning) theo khung của Pearl \cite{pearl2018why} phân biệt ba tầng:

\begin{enumerate}[leftmargin=1.5cm, itemsep=4pt]
  \item \textbf{Tương quan (Association):} ``Những ngày mưa, lưu lượng tại nút A thấp hơn.''
  \item \textbf{Can thiệp (Intervention):} ``Nếu tôi đóng đường X [do-calculus], lưu lượng ở Y sẽ thay đổi thế nào?''
  \item \textbf{Phản thực tế (Counterfactual):} ``Nếu thay vào đó tôi đóng đường Z (thay vì X), hậu quả sẽ tệ hơn hay tốt hơn?''
\end{enumerate}

\subsection{CF-VLA và Cơ chế Lấy cảm hứng trong STWI}

CF-VLA là mô hình Vision--Language--Action end-to-end cho lái xe tự động. STWI không triển khai kiến trúc đó; dự án chỉ lấy cảm hứng từ ý tưởng self-reflection để xây dựng Counterfactual Safety Loop trên dữ liệu, mô phỏng và citations.

\begin{tcolorbox}[dangerbox]
\textbf{Vai trò của Safety Loop:} Kiểm tra tác động bậc hai trước khi phát hành khuyến nghị. Khi không hội tụ, OOD hoặc thiếu evidence, workflow fail closed và chuyển operator review.
\end{tcolorbox}

\section{Độ Đo Đánh giá Hệ thống}

Các độ đo chính được sử dụng xuyên suốt trong báo cáo:

\begin{align}
  \text{RMSE} &= \sqrt{\frac{1}{n}\sum_{i=1}^n (\hat{y}_i - y_i)^2} \label{eq:rmse} \\
  \text{MAE}  &= \frac{1}{n}\sum_{i=1}^n |\hat{y}_i - y_i| \label{eq:mae} \\
  \text{F1}   &= 2 \cdot \frac{\text{Precision} \times \text{Recall}}{\text{Precision} + \text{Recall}} \label{eq:f1} \\
  \text{Recall@}k &= \frac{|\text{Rel}_k \cap \text{Ret}_k|}{|\text{Rel}|} \label{eq:recall_k}
\end{align}

Trong đó $\hat{y}_i$ là giá trị dự báo, $y_i$ là giá trị thực tế, $\text{Rel}_k$ là tập tài liệu liên quan trong top-$k$ kết quả truy xuất, $\text{Ret}_k$ là tập tài liệu được truy xuất.

\section{Công trình Liên quan (Related Work)}

Phần này tổng hợp và so sánh các công trình nghiên cứu liên quan trực tiếp đến bốn trụ cột kỹ thuật chính của hệ thống STWI: dự báo giao thông dựa trên đồ thị, mô hình thay thế cho mô phỏng, truy xuất tri thức pháp lý, và hệ thống hỗ trợ ra quyết định giao thông.

\subsection{Dự báo Giao thông dựa trên Đồ thị}

Các mô hình dự báo giao thông hiện đại khai thác cấu trúc không gian--thời gian (spatio-temporal) của mạng lưới đường phố thông qua mạng nơ-ron đồ thị. Hai công trình tiêu biểu trong hướng nghiên cứu này là STGCN và DCRNN.

\textbf{STGCN} \cite{yu2018stgcn} đề xuất kiến trúc kết hợp phép tích chập đồ thị Chebyshev trong miền phổ với tích chập nhân quả 1D (causal convolution) theo chiều thời gian. Mô hình sử dụng cấu trúc ``sandwich'' ST-Conv Block gồm hai lớp tích chập thời gian bao quanh một lớp tích chập đồ thị, cho phép trích xuất đồng thời các đặc trưng không gian và thời gian. Ưu điểm chính của STGCN là tốc độ huấn luyện nhanh nhờ tránh được cơ chế hồi quy (recurrence), nhưng khả năng mô hình hóa phụ thuộc dài hạn bị giới hạn bởi kích thước cửa sổ tích chập.

\textbf{DCRNN} \cite{li2018diffusion} mô hình hóa quá trình lan truyền giao thông như phép khuếch tán (diffusion process) trên đồ thị có hướng, sử dụng hai ma trận chuyển tiếp (forward và backward) để nắm bắt hướng của dòng giao thông. Phần thời gian được xử lý bằng GRU (Gated Recurrent Unit) theo cơ chế encoder--decoder. DCRNN đạt hiệu quả cao trên các tập dữ liệu chuẩn (METR-LA, PEMS-BAY) nhưng có chi phí tính toán lớn do cơ chế hồi quy tuần tự.

\textbf{Cách tiếp cận của STWI:} Trong phiên bản MVP, STWI sử dụng kiến trúc \textbf{GCN + Stacked LSTM} theo cách xếp chồng đơn giản: các lớp GCN (Mục~2.1) trích xuất đặc trưng không gian, sau đó Stacked LSTM (Mục~2.2) mô hình hóa phụ thuộc thời gian. So với STGCN và DCRNN, kiến trúc này có ưu điểm dễ triển khai, dễ gỡ lỗi và phù hợp với quy mô dữ liệu ban đầu của TP.HCM. Lộ trình phát triển dự kiến bổ sung \textbf{cơ chế attention} (spatial attention và temporal attention) ở các phiên bản sau để cải thiện khả năng nắm bắt các phụ thuộc không gian--thời gian phức tạp hơn.

\subsection{Mô hình Thay thế cho Mô phỏng Giao thông}

Mô phỏng vi mô (SUMO, Vissim) cung cấp kết quả chính xác nhưng chi phí tính toán quá lớn cho ứng dụng thời gian thực. Nghiên cứu gần đây tập trung xây dựng các mô hình thay thế nơ-ron (neural surrogate) để xấp xỉ kết quả mô phỏng với tốc độ nhanh hơn nhiều bậc.

\textbf{Wen et al.} \cite{wen2023surrogate} tiến hành nghiên cứu so sánh toàn diện các kiến trúc neural surrogate cho SUMO, bao gồm MLP, CNN, và các biến thể có cơ chế attention. Nghiên cứu chỉ ra rằng việc kết hợp nhiều kiến trúc (ensemble) cải thiện đáng kể độ chính xác so với mô hình đơn lẻ, đặc biệt khi xử lý các kịch bản giao thông đa dạng. Tuy nhiên, nghiên cứu chưa đề cập đến cơ chế phát hiện các mẫu đầu vào nằm ngoài phân phối huấn luyện (Out-of-Distribution).

\textbf{Cách tiếp cận của STWI:} Hệ thống STWI sử dụng \textbf{heterogeneous ensemble} gồm ba kiến trúc khác nhau (MLP, CNN-1D, Transformer nhẹ) để tăng đa dạng dự báo và cải thiện ước lượng độ bất định (uncertainty) theo phương pháp Deep Ensemble \cite{lakshminarayanan2017ade}. Điểm khác biệt chính so với các công trình liên quan là STWI bổ sung \textbf{OOD Detection Gate}: khi phương sai ensemble vượt ngưỡng hoặc đầu vào được phân loại ngoài phân phối, hệ thống tự động chuyển sang chế độ fail-closed, từ chối đưa ra dự báo và chuyển giao cho operator review. Cơ chế này đảm bảo tính an toàn trong các tình huống mà mô hình chưa gặp trong quá trình huấn luyện.

\subsection{RAG và Truy xuất Tri thức Pháp lý}

Truy xuất Tăng cường Phát sinh (RAG) là kỹ thuật then chốt cho việc kết hợp LLM với nguồn tri thức bên ngoài, đặc biệt quan trọng trong bối cảnh pháp lý nơi tính chính xác và trích dẫn nguồn là yêu cầu bắt buộc.

\textbf{RAG cổ điển} \cite{lewis2020rag} kết hợp bộ truy xuất dày đặc (dense retriever) với bộ sinh văn bản (generator) theo kiến trúc hai giai đoạn: truy xuất đoạn văn bản liên quan từ kho tri thức, sau đó đưa vào ngữ cảnh của LLM để sinh câu trả lời. Phương pháp này hiệu quả cho văn bản phi cấu trúc nhưng gặp khó khăn khi truy vấn dữ liệu có cấu trúc (bảng biểu, số liệu mô phỏng).

\textbf{XiYanSQL} \cite{guo2023xiyansql} giải quyết bài toán Text-to-SQL bằng kiến trúc multi-generator ensemble, trong đó nhiều bộ sinh SQL hoạt động song song và kết quả được chọn lọc qua bước xác minh. Phương pháp này cho phép truy vấn dữ liệu có cấu trúc từ ngôn ngữ tự nhiên nhưng yêu cầu LLM sinh trực tiếp câu lệnh SQL --- tiềm ẩn rủi ro SQL injection và khó kiểm soát phạm vi truy vấn.

\textbf{Cách tiếp cận của STWI:} Thay vì để LLM sinh SQL thô, STWI áp dụng phương pháp \textbf{Schema-Level RAG}: LLM tạo ra đối tượng \textit{SimulationQuery} có kiểu dữ liệu rõ ràng (typed query object), sau đó phía server chuyển đổi thành câu lệnh SQL tham số hóa (parameterized SQL). Phương pháp này loại bỏ hoàn toàn rủi ro SQL injection và giới hạn phạm vi truy vấn trong schema đã định nghĩa trước. Ngoài ra, STWI bổ sung cơ chế \textbf{citation validation với kiểm tra ngày hiệu lực} (effective-date checking): mỗi trích dẫn pháp lý được xác minh không chỉ về nội dung mà còn về thời hạn hiệu lực của văn bản pháp quy, đảm bảo khuyến nghị luôn dựa trên quy định hiện hành \cite{gtvt2023law}.

\subsection{Hệ thống Hỗ trợ Ra quyết định Giao thông}

Hệ thống Giao thông Thông minh (ITS -- Intelligent Transport Systems) truyền thống dựa trên luật cố định (rule-based): các ngưỡng lưu lượng, tốc độ được cấu hình tĩnh, khi vượt ngưỡng thì kích hoạt cảnh báo hoặc điều chỉnh đèn tín hiệu. Ưu điểm là minh bạch và dễ kiểm toán, nhưng không thích ứng được với các tình huống mới và thiếu khả năng suy luận ``what-if'' (điều gì sẽ xảy ra nếu thay đổi một biến số).

Các hệ thống AI-based gần đây tích hợp mô hình học sâu để dự báo và tối ưu hóa tín hiệu giao thông, nhưng thường tập trung vào một khía cạnh duy nhất (chỉ dự báo, hoặc chỉ tối ưu đèn) mà thiếu khung tích hợp toàn diện từ dữ liệu đến ra quyết định.

\textbf{Cách tiếp cận của STWI:} Hệ thống STWI kết hợp bốn thành phần trong một khung thống nhất: (i)~kiến trúc \textbf{đa tác tử} \cite{wooldridge2009multiagent} cho phép phân rã nhiệm vụ phức tạp thành các tác tử chuyên biệt (Simulation, Legal, Evaluation); (ii)~\textbf{Counterfactual Safety Loop} \cite{pearl2018why} kiểm tra tác động bậc hai trước khi phát hành khuyến nghị; (iii)~\textbf{Legal RAG} \cite{lewis2020rag} đảm bảo mọi đề xuất đều có cơ sở pháp lý kèm trích dẫn; và (iv)~thiết kế \textbf{fail-closed} -- khi bất kỳ thành phần nào không hội tụ hoặc phát hiện OOD, toàn bộ workflow dừng lại và chuyển sang operator review. Sự kết hợp đồng thời bốn yếu tố này là đặc trưng riêng của STWI so với các hệ thống hiện có.

\subsection{Bảng So sánh Tổng hợp}

Bảng~\ref{tab:related_work_comparison} tổng hợp so sánh STWI với các hệ thống và công trình liên quan trên các chiều chính.

\begin{table}[H]
\centering
\caption{So sánh STWI với các công trình và hệ thống liên quan}
\label{tab:related_work_comparison}
\renewcommand{\arraystretch}{1.4}
\small
\begin{tabular}{p{2.6cm} p{2.8cm} p{2.2cm} p{2.5cm} p{2.5cm}}
\toprule
\textbf{Hệ thống / Công trình} & \textbf{Kiến trúc} & \textbf{Thời gian thực} & \textbf{Cơ chế An toàn} & \textbf{Tuân thủ Pháp lý} \\
\midrule
STGCN \cite{yu2018stgcn} & ST-Conv Block (Chebyshev + causal conv) & Có (dự báo) & Không & Không \\
\rowcolor{rowgray}
DCRNN \cite{li2018diffusion} & Diffusion conv + GRU encoder--decoder & Có (dự báo) & Không & Không \\
Wen et al. \cite{wen2023surrogate} & Neural surrogate cho SUMO & Có (surrogate) & Không & Không \\
\rowcolor{rowgray}
RAG cổ điển \cite{lewis2020rag} & Dense retriever + generator & Có (truy vấn) & Không & Một phần (trích dẫn) \\
XiYanSQL \cite{guo2023xiyansql} & Multi-generator Text-to-SQL & Có (truy vấn) & Không & Không \\
\rowcolor{rowgray}
ITS truyền thống & Rule-based, ngưỡng cố định & Có & Giới hạn (luật cứng) & Tĩnh (cấu hình thủ công) \\
\textbf{STWI (MVP)} & \textbf{Multi-agent + GCN-LSTM + Surrogate + RAG} & \textbf{Có (end-to-end)} & \textbf{OOD gate + CF Safety Loop + fail-closed} & \textbf{Legal RAG + effective-date check} \\
\bottomrule
\end{tabular}
\end{table}

\begin{tcolorbox}[notebox]
\textbf{Nhận xét:} Không có công trình nào trong danh sách trên kết hợp đồng thời cả bốn yếu tố: dự báo spatio-temporal, mô phỏng what-if thời gian thực, truy xuất tri thức pháp lý có kiểm tra hiệu lực, và vòng lặp an toàn phản thực tế. STWI hướng đến việc tích hợp các thành phần này trong một kiến trúc thống nhất, trong đó mỗi thành phần được thiết kế fail-closed để đảm bảo an toàn tổng thể của hệ thống.
\end{tcolorbox}
